\documentclass[11pt]{amsart}
\usepackage[T1]{fontenc}
\usepackage{lmodern}
\usepackage[a4paper,margin=30mm]{geometry}
\usepackage{amsmath,amssymb,mathtools}
\usepackage{microtype}
\usepackage[colorlinks=true,linkcolor=blue,citecolor=blue,urlcolor=blue]{hyperref}
\hypersetup{pdftitle={Asymptotic matching and initial-set dependence in the precision cost of contracting control systems}}

\newtheorem{theorem}{Theorem}[section]
\newtheorem{proposition}[theorem]{Proposition}
\newtheorem{lemma}[theorem]{Lemma}
\newtheorem{corollary}[theorem]{Corollary}
\theoremstyle{definition}
\newtheorem{definition}[theorem]{Definition}
\theoremstyle{remark}
\newtheorem{remark}[theorem]{Remark}
\newcommand{\R}{\mathbb R}
\newcommand{\N}{\mathbb N}
\newcommand{\area}{\operatorname{area}}
\newcommand{\dist}{\operatorname{dist}}
\newcommand{\diam}{\operatorname{diam}}
\newcommand{\ee}{\varepsilon}
\newcommand{\doi}[1]{\href{https://doi.org/#1}{doi:#1}}
\numberwithin{equation}{section}



\title[Asymptotic matching and precision cost]{Asymptotic matching and initial-set dependence\\
in the precision cost of contracting control systems}
\author{Author information pending}
\date{}
\subjclass[2020]{Primary 93C55, 37B40; Secondary 93D15, 37C05}
\keywords{Outer invariance complexity, precision cost, controlled contraction,
asymptotic matching, feasible geometry, initial-set dependence}

\begin{document}
\begin{abstract}
We study the minimum number of open-loop controls required to keep
every initial state within a Euclidean tolerance of a constraint set
at all times of a finite horizon. For an explicit class of smooth
planar contracting systems, feasible branch lengths and radial
contraction may depend independently on the state. Matching at the
common attracting origin, together with directly checkable small
derivative bounds, yields a common-cost boundary $\gamma_*=\beta h_a$,
where $h_a=-\sum_i a_i\log a_i$ and
$\delta_n=\exp(-\gamma n+o(n))$. At and below this boundary, every
fixed positive-area compact initial set has rate $\gamma/\beta$.
Above it, one compact set of arbitrarily nearly full area has rate
$h_a$ for all finite higher exponents and their precision sequences,
whereas the full constraint set has a strictly larger lower rate.
The proof compares arbitrary approximate controls with stopped
physical witnesses at a loss of at most $1+Bn$. It does not require
an instantaneous power identity between branch lengths and radial
factors. We also give instances excluding coordinate reduction to
specified exactly matched models. A separate radius-autonomous
mismatch theorem and a smooth collapse example delimit the matching
condition and the role of global controlled geometry.
\end{abstract}
\maketitle

\section{Introduction}

Invariance entropy measures the control information required to keep
an initial set inside a controlled invariant set; see
Kawan \cite{Kawan2013}. When all controls contract toward a point,
outer invariance complexity can be bounded at every fixed positive
tolerance. An exponentially decreasing tolerance can nevertheless
require exponentially many controls. We ask when this precision
cost is shared by all fixed positive-area compact initial sets,
and when the geometry of the initial set changes its rate.

The catalogue used here assigns an open-loop input to each actual
initial state. One tolerance $\delta_n$ applies at every integer time
from $0$ through $n$. A trajectory may leave and return to the
constraint set within that tolerance. For this task, neither an exact
coding nor convergence of all trajectories alone determines the
optimal catalogue. An almost full-area set may exclude the rare
states which determine the full-set cost.

Our main result treats two complementary feasible branches with
limiting proportions $a_0,a_1$ and radial first derivatives
$a_0^\beta,a_1^\beta$. The branch proportions and radial factors are
independent smooth functions of the state. A smallness condition on
these functions gives complete feasible intervals, even though
states with the same initial radius and control word may have
different subsequent radial trajectories. Their cumulative angular
and radial errors are bounded separately. Thus instantaneous
matching is unnecessary: the boundary $\beta h_a$ survives this
specified class of perturbations.

The controlled geometric step concerns arbitrary approximate inputs.
A logarithmic slope cone governs the moving feasible graphs. At a
first wrong branch, the physical excess outside the constraint yields
a one-sided band in the actual initial fibre. Complete stopped
intervals provide separated witnesses, and one input covers at most
$1+Bn$ of them. A global bound toward the origin supplies the entire
tail after an early stop. Only after these estimates have been proved
do classical stopping and type counts give the precision rates.

Theorem~\ref{thm:main} proves the common-cost boundary. A separate
radius-autonomous family, Theorem~\ref{thm:mismatch}, shows that
unequal radial matching powers cause initial-set dependence at every
sufficiently small positive exponent. This necessity statement stays
within that family. Section~\ref{sec:conjugacy} gives instances of
the main theorem which cannot be reduced to the specified previous
forms by a simultaneous change of coordinates preserving the
constraint. Section~\ref{sec:local} shows that agreement near the
origin and all-input convergence, without the global geometric
conditions, do not imply the boundary.

\subsection{Relation to previous work}

The finite-horizon count below is the outer invariance complexity
of Chen and Zhong \cite[Section 2.2]{CZ2024}, with $n+1$ constrained
times instead of $n$. Their fixed-tolerance complexity and
equi-invariability results provide the underlying catalogue framework.
Exponential and practical stabilization entropies
\cite{Colonius2012,CH2021} address related information costs for
different stabilization tasks. We do not introduce a new catalogue
definition.

Measure-theoretic invariance entropies and variational principles
are developed in \cite{WHS2019,NWH2022}. The entropy dimensions of
Chen, Huang and Zhong \cite{CHZ2026} use invariant-partition
cylinders and weights $\exp[-s(N\tau)^\alpha]$, with a separate
measure-covering tolerance. Their cylinder and variational estimates
do not by themselves control the Euclidean first-deviation bands of
arbitrary inputs at $\delta_n$. On our connected $Q$, a finite
clopen invariant partition would have only one nonempty cell, while
neither control preserves all of $Q$. The main clopen variational
theorems therefore do not directly apply. Measurable-cylinder
estimates can still be useful, but require the physical comparison
proved here. Ordinary area need not be invariant; its use is not
a new measure-theoretic framework.

Graph-transform and volume-distortion methods are existing tools.
Kawan and Da Silva \cite[Appendix]{KD2018} prove volume estimates
near a partially hyperbolic all-time lift. For our contracting
systems, every bounded bi-infinite trajectory is the origin:
$\|x_0\|_*\le L^k\sup_j\|x_j\|_*$ for all $k$.
The physical angular coordinate is singular there. Those local
volume estimates alone do not give the complete feasible splits
or the arbitrary-input comparison as the tolerance decreases with
the horizon.

Prefix stopping, product weights, types, and typical sets are
classical \cite{Hutchinson1981,Csiszar1998}. Yang, Chen, Yang and
Zhou \cite{YCYZ2025} develop cost-threshold pressure and
Bowen-equation results on a given invariant partition. Their
spanning/separated comparison preserves the prescribed address;
it does not compare that address with every approximately feasible
physical input. The additional step needed here is the band
estimate and stopped-witness comparison of
Section~\ref{sec:catalogue}. These comparisons with
\cite{KD2018,YCYZ2025} use the graph-transform appendix and
the author versions arXiv:1711.01181v2 and arXiv:2309.01628v3,
respectively. Our proofs of the required estimates are self-contained.

\section{Counting problem and main results}

Let
\[
Q=\{(s,z)\in\R^2:0\le s\le1,\ |z|\le s\},\qquad \area(Q)=1.
\]
For two maps $F_0,F_1:\R^2\to\R^2$ and
$u=(u_1,u_2,\ldots)\in\{0,1\}^{\N}$, write
$F_u^0x=x$ and $F_u^kx=F_{u_k}\circ\cdots\circ F_{u_1}x$.
All binary inputs are allowed. Distances are Euclidean, and
$Q^\delta=\{x:\dist(x,Q)<\delta\}$.

\begin{definition}\label{def:catalogue}
For a nonempty compact $K\subset Q$, let $r_F(n,\delta;K)$ be
the least size of a finite catalogue $\mathcal U\subset\{0,1\}^{\N}$
such that
\begin{equation}\label{eq:catalogue}
\forall x\in K\ \exists u\in\mathcal U\
\forall k\in\{0,\ldots,n\}:\quad F_u^kx\in Q^\delta.
\end{equation}
The input assignment may depend on the actual state. A trajectory may
leave and reenter $Q$, subject to every constraint in
\eqref{eq:catalogue}. Only the first $n$ letters matter; infinite
inputs make tail constructions unambiguous. Exact controlled
feasibility will prove that this minimum is finite.
\end{definition}

A precision sequence is positive, tends to zero, and satisfies
$-\log\delta_n/n\to\gamma$ for a finite $\gamma\ge0$.
It need not be monotone. All logarithms are natural. Fix
\begin{equation}\label{eq:parameters}
a_0,a_1\in(0,1),\quad a_0+a_1=1,\quad a_0\ne a_1,\quad\beta>1,
\end{equation}
and put $a_* =\min_i a_i$, $a^+=\max_i a_i$, $d=-\log a^+$, and
$h=-\sum_i a_i\log a_i$.
For later counting, set
\[
H_b(v)=-v\log v-(1-v)\log(1-v),\qquad
\chi_a(v)=-v\log a_0-(1-v)\log a_1,
\]
with $0\log0=0$. Thus $H_b(a_0)=\chi_a(a_0)=h$.

\subsection{Independent branch and radial perturbations}

Define the sufficient threshold
\begin{equation}\label{eq:smallthreshold}
\nu_0=\min\left\{\frac{a_*}{1000},
      \frac{(\beta-1)d}{4(1+2/a_*)}\right\}>0.
\end{equation}
Fix $0<\nu<\nu_0$ and three independent functions
$\eta,\xi_0,\xi_1\in C_c^\infty(B(0,2))$, each zero at the origin,
such that
\begin{equation}\label{eq:functionbounds}
\max_{f\in\{\eta,\xi_0,\xi_1\}}
\max_{|\alpha|\le2}\|\partial^\alpha f\|_\infty\le\nu.
\end{equation}
Here $B(0,2)$ is the fixed Euclidean open ball in $\R^2$.
At each input state $(s,z)$ put
\begin{equation}\label{eq:independentfactors}
p_0=a_0+\eta,\quad p_1=a_1-\eta,
\qquad R_i=a_i^\beta e^{\xi_i},
\end{equation}
and define on the entire plane
\begin{align}
F_0(s,z)&=\left(sR_0,\frac{R_0}{p_0}(z+p_1s)\right),\nonumber\\
F_1(s,z)&=\left(sR_1,\frac{R_1}{p_1}(z-p_0s)\right).
\label{eq:F}
\end{align}
There is no imposed relation $R_i=p_i^\beta$.
The matching at the origin concerns the matrices
\begin{equation}\label{eq:firstjets}
\mathsf A_0=\begin{pmatrix}a_0^\beta&0\\a_0^{\beta-1}a_1&a_0^{\beta-1}\end{pmatrix},
\qquad
\mathsf A_1=\begin{pmatrix}a_1^\beta&0\\-a_1^{\beta-1}a_0&a_1^{\beta-1}\end{pmatrix}.
\end{equation}

\begin{theorem}[Common-cost boundary under asymptotic matching]\label{thm:main}
For every triple satisfying \eqref{eq:smallthreshold}--\eqref{eq:functionbounds},
the maps \eqref{eq:F} are global $C^\infty$ diffeomorphisms,
fix the origin with $DF_i(0)=\mathsf A_i$, drive bounded initial
sets uniformly exponentially to the origin under all inputs,
and make $Q$ controlled invariant.
For every fixed positive-area compact $K\subset Q$ and every
precision sequence of finite exponent $\gamma\ge0$,
\begin{equation}\label{eq:universal}
\min\{h,\gamma/\beta\}\le
\liminf_n\frac{\log r_F(n,\delta_n;K)}n
\le\limsup_n\frac{\log r_F(n,\delta_n;K)}n
\le\min\{\log2,\gamma/\beta\}.
\end{equation}
In particular, for $0\le\gamma\le\beta h$ the limit exists
and equals $\gamma/\beta$.

For every $\zeta\in(0,1)$ there is a fixed compact $K_\zeta\subset Q$
with $\area(K_\zeta)>1-\zeta$ such that, simultaneously for every
finite $\gamma>\beta h$ and every precision sequence of that exponent,
\begin{equation}\label{eq:high}
\lim_n\frac{\log r_F(n,\delta_n;K_\zeta)}n=h
<\liminf_n\frac{\log r_F(n,\delta_n;Q)}n.
\end{equation}
This set may depend on the fixed system and $\zeta$, but is chosen
before the exponent, precision sequence, and horizon.
\end{theorem}

The threshold concerns shared rates across all positive-area compact
sets. It does not classify the possible rates above the threshold
or give an exact high-precision spectrum or limit for $Q$.
The bound \eqref{eq:smallthreshold} is sufficient and is not claimed
to be the largest allowable perturbation.

\subsection{A different family: radius-autonomous mismatch}
Keep $a_0,a_1$ as above, but now let $\beta_0,\beta_1>1$ and
$\beta_{\max}=\max_i\beta_i$. For
\begin{equation}\label{eq:radialepsilon}
0\le\ee<\frac{a_*}{4(\beta_{\max}+1)},\qquad
q_0(s)=a_0+\frac{\ee s}{1+s^2},\quad
q_1(s)=a_1-\frac{\ee s}{1+s^2},
\end{equation}
define
\begin{align}
G_0(s,z)&=(sq_0(s)^{\beta_0},q_0(s)^{\beta_0-1}(z+q_1(s)s)),\nonumber\\
G_1(s,z)&=(sq_1(s)^{\beta_1},q_1(s)^{\beta_1-1}(z-q_0(s)s)).\label{eq:G}
\end{align}
Set
\begin{equation}\label{eq:mismatchconstants}
\rho_i=a_i^{\beta_i},\quad
\chi=-\sum_i a_i\log\rho_i,\quad
\gamma_0=-\log\max_i\rho_i,\quad
\rho_0^D+\rho_1^D=1\quad(D>0).
\end{equation}

\begin{theorem}[Low-precision separation under radial mismatch]\label{thm:mismatch}
Suppose $\beta_0\ne\beta_1$. The maps \eqref{eq:G} are global
smooth diffeomorphisms with a common attracting origin under all
inputs, and $Q$ is controlled invariant. For every $\zeta\in(0,1)$
there is one compact $K_\zeta\subset Q$ of area greater than $1-\zeta$
such that, simultaneously for every $0<\gamma<\gamma_0$ and every
precision sequence of exponent $\gamma$,
\begin{equation}\label{eq:mismatchrates}
\lim_n\frac{\log r_G(n,\delta_n;Q)}n=\gamma D,
\qquad
\lim_n\frac{\log r_G(n,\delta_n;K_\zeta)}n=\gamma h/\chi,
\qquad D>h/\chi.
\end{equation}
Every fixed positive-area compact $K\subset Q$ also has lower rate
at least $\gamma h/\chi$ in this range.
\end{theorem}

Thus a fixed system can exhibit strict initial-set dependence at
every sufficiently small positive exponent. The theorem is restricted
to \eqref{eq:G}; it makes no claim about unequal powers in the
independent-perturbation family. The precise matching criterion in this
radius-autonomous class is given in Corollary~\ref{cor:criterion}.


\section{Complete feasible graphs and independent distortion}
\label{sec:geometry}

Throughout this section the system is \eqref{eq:F}. Set
\[
p_-=a_*-\nu,\quad p_+=1-p_-,\quad r=e^\nu(a^+)^\beta,
\quad t=\frac{e^\nu(a^+)^{\beta-1}}{1-\nu/a_*},
\quad\sigma_0=-1,\quad\sigma_1=1.
\]
The following bounds apply on the whole plane:
$p_-\le p_i\le p_+$, $R_i\le r$, and $R_i/p_i\le t$.

\begin{lemma}[Global dynamics and one-step feasibility]\label{lem:dynamics}
The dynamical assertions of Theorem~\ref{thm:main} hold.
For $s>0$, write $y=z/s$, $P_i(s,y)=p_i(s,sy)$, and
$W_i(s,y)=\xi_i(s,sy)$.
There is a unique $\theta(s)\in(-1,1)$ satisfying
\begin{equation}\label{eq:divider}
\theta(s)=2P_0(s,\theta(s))-1.
\end{equation}
The complete one-step feasible domains are $[-1,\theta(s)]$
for control $0$ and $[\theta(s),1]$ for control $1$.
Choosing $0$ at the shared endpoint defines a natural exact code.
\end{lemma}

\begin{proof}
Since $\nu<a_*/2$, $-\log(1-\nu/a_*)\le2\nu/a_*$. Hence
\begin{equation}\label{eq:contractionfactors}
\log t\le-(\beta-1)d+(1+2/a_*)\nu
<-\tfrac34(\beta-1)d<0,
\quad \log r<-\tfrac34\beta d<0.
\end{equation}
Take
\begin{equation}\label{eq:origincontraction}
M=\frac{1+t}{1-t},\quad
\|x\|_* =\max(M|s|,|z|),\quad
L=\max\{r,2t/(1+t)\}<1.
\end{equation}
The second coordinate of \eqref{eq:F} has coefficient at most $t$
in absolute value on both $z$ and $s$. Thus, on the entire plane,
\begin{equation}\label{eq:globaloriginbound}
\|F_i x\|_*\le L\|x\|_*.
\end{equation}
This bounds contraction toward the origin; it does not assert
pairwise contraction. It also controls tails outside $Q$.

For global invertibility use the linear shear
$J_i(s,z)=(s,z-\sigma_i s)=(s,v)$, whose norm and inverse norm
are less than $2$. Input and output shears of \eqref{eq:F} give
$(S,V)=(sR_i,vR_i/p_i)$, because $p_0+p_1=1$.
If $g_i=\mathsf A_i^{-1}F_i-\operatorname{id}$, then
\[
J_i g_i(x)=(s\alpha_i(x),v\omega_i(x)),\qquad
\alpha_i=e^{\xi_i}-1,\quad \omega_i=e^{\xi_i}a_i/p_i-1.
\]
This $g_i$ vanishes outside $B(0,2)$. On its support
$|s|\le2$, $|v|\le2\sqrt2$, and differentiation gives
\begin{align*}
|\alpha_i|&\le e^\nu\nu,&
\|\nabla\alpha_i\|_2&\le e^\nu\sqrt2\nu,\\
|\omega_i|&\le e^\nu\nu(1+1/p_-),&
\|\nabla\omega_i\|_2&\le2e^\nu\sqrt2\nu(1+1/p_-).
\end{align*}
For the last bound use $a_i/p_i\le2$.
The two row-gradient estimates, followed by the inverse shear, yield
\begin{equation}\label{eq:inversebound}
\|Dg_i\|_2\le
2e^\nu\nu\bigl[(1+2\sqrt2)+(8+\sqrt2)(1+1/p_-)\bigr]
<200\nu/a_*<1/5.
\end{equation}
The coarse bound follows from $e^\nu<3$, $p_->a_*/2$, and
$a_*\le1/2$. It holds globally, since $g_i$ has compact support.
For each target $y\in\R^2$, the equation $F_i x=y$ is equivalent
to $x=\mathsf A_i^{-1}y-g_i(x)$, a contraction equation on the
complete Euclidean plane. It has a unique solution.
Also $DF_i=\mathsf A_i(I+Dg_i)$ is nonsingular everywhere.
The local inverse theorem now gives a smooth global inverse.
Positivity of $p_i$ makes \eqref{eq:F} smooth, and the zero values
of the perturbations give $DF_i(0)=\mathsf A_i$; their first
derivatives need not vanish.

In the calculation coordinates the actual map is
\begin{equation}\label{eq:angularmap}
S=sa_i^\beta e^{W_i(s,y)},\qquad
Y=\sigma_i+(y-\sigma_i)/P_i(s,y).
\end{equation}
The physical set and metric are unchanged. On $0<s\le1$, $|y|\le1$,
the value at zero and the derivative bounds imply
\begin{equation}\label{eq:partialbounds}
\begin{split}
|P_i-a_i|,\ |W_i|&\le2\nu s,\\
|(P_i)_s|,\ |(W_i)_s|&\le2\nu,
\qquad |(P_i)_y|,\ |(W_i)_y|\le\nu s.
\end{split}
\end{equation}
The function $y-2P_0(s,y)+1$ has derivative at least
$1-2\nu>0$ and opposite signs at $-1,1$, proving the unique
divider. The conditions $-1\le Y\le1$ give exactly the asserted
domains, each mapped angularly onto $[-1,1]$. Since $S\le rs<s$,
there is no additional cap constraint. The vertex is fixed.
These facts prove controlled invariance and the Borel exact code.
\end{proof}

\begin{lemma}[Moving graph cone and complete intervals]\label{lem:graphs}
Fix an actual initial radius $0<s_0\le1$. For every finite word $w$,
the initial angles feasible at all times through $|w|$ form a
nondegenerate closed interval $I_w(s_0)$. Its terminal angular map
is increasing onto $[-1,1]$, and its terminal radial graph $s=S(y)$
satisfies
\begin{equation}\label{eq:cone}
|S'(y)/S(y)|\le1.
\end{equation}
Every complete prefix-free word family partitions the initial fibre,
with disjoint interiors and all shared endpoints retained.
\end{lemma}

\begin{proof}
Suppose a parent graph obeys \eqref{eq:cone}. Put $k=S'/S$ and
\[
A_i=(P_i)_sS'+(P_i)_y,\qquad
B_i=(W_i)_sS'+(W_i)_y.
\]
Equation \eqref{eq:partialbounds} gives $|A_i|,|B_i|\le3\nu S$.
Differentiating along this actual graph gives
\begin{equation}\label{eq:graphupdate}
q_i=1-(y-\sigma_i)A_i/P_i,\quad
\frac{dY}{dy}=q_i/P_i,\quad
k_{\rm new}=P_i(k+B_i)/q_i.
\end{equation}
The radial term $B_i$ is independent of $A_i$.
In particular $|q_i-1|\le6\nu S/P_i$ and $q_i>0$.
Since $P_i(1-P_i)\ge p_-(1-p_-)\ge a_*/4$ and
$(3P_i^2+6)\nu S\le9\nu<a_*/4$, we have
\begin{equation}\label{eq:coneinvariance}
|k_{\rm new}|\le
\frac{P_i(1+3\nu S)}{1-6\nu S/P_i}<1.
\end{equation}
On this graph $\Psi(y)=y-2P_0(S(y),y)+1$ has derivative at least
$1-6\nu S>0$ and opposite endpoint signs. It divides the parent
once. Each feasible side is mapped increasingly onto the whole
$[-1,1]$, and \eqref{eq:coneinvariance} applies to its output graph.

The original fibre has $S\equiv s_0$ and $k=0$.
Induction proves the assertions at every time. Splitting each
parent interval proves the partition assertion. The positive
derivatives also give smooth dependence of finite-word endpoints
on $s_0>0$ by the implicit function theorem. Closed endpoints
and every intermediate constraint have been retained.
\end{proof}

For a word $w=w_1\cdots w_\ell$, set
$\lambda_w=\prod_{j=1}^\ell a_{w_j}$, with empty product $1$, and write
\[
\mathcal A_w=\{(s_0,s_0y):0<s_0\le1,\ y\in I_w(s_0)\}\cup\{0\}
\]
for its complete physical feasible strip.

\begin{lemma}[Separate radial and angular comparisons]\label{lem:products}
Fixed constants $C_R,C_A\ge1$ satisfy, for every complete feasible
word and every point of its initial interval,
\begin{align}
C_R^{-1}s_0\lambda_w^\beta&\le s_\ell\le C_Rs_0\lambda_w^\beta,
\nonumber\\
C_A^{-1}\lambda_w^{-1}&\le\partial_{y_0}Y_{w,s_0}
\le C_A\lambda_w^{-1}.\label{eq:products}
\end{align}
Consequently
\begin{equation}\label{eq:intervalsize}
2C_A^{-1}\lambda_w\le|I_w(s_0)|\le2C_A\lambda_w,
\end{equation}
and
\begin{equation}\label{eq:cylinderarea}
C_A^{-1}\lambda_w\le\area(\mathcal A_w)\le C_A\lambda_w.
\end{equation}
Finite-word boundary curves have zero area and remain in the covers.
\end{lemma}

\begin{proof}
Along a complete feasible word $s_j\le r^js_0$.
Set
\begin{gather*}
B_P=\frac{2\nu}{p_-(1-r)},\quad B_R=\frac{2\nu}{1-r},
\quad d_0=\frac{6\nu}{p_-}<1,\\
B_q=\frac{6\nu}{p_-(1-d_0)(1-r)},\quad
C_R=e^{B_R},\quad C_A=e^{B_P+B_q}.
\end{gather*}
Equations \eqref{eq:partialbounds} and \eqref{eq:graphupdate} imply
\begin{equation}\label{eq:logerrors}
\begin{split}
\left|\sum_{j<\ell}\log\frac{p_{w_{j+1}}(x_j)}{a_{w_{j+1}}}\right|
&\le B_P,\qquad
\left|\sum_{j<\ell}\xi_{w_{j+1}}(x_j)\right|\le B_R,\\
\sum_{j<\ell}|\log q_{w_{j+1}}|&\le B_q.
\end{split}
\end{equation}
For the last estimate use $|\log q|\le|q-1|/(1-d_0)$.
The identities
\[
s_\ell=s_0\lambda_w^\beta
 \exp\!\sum_{j<\ell}\xi_{w_{j+1}}(x_j),\qquad
\partial_{y_0}Y_{w,s_0}=\prod_{j<\ell}q_{w_{j+1}}/P_{w_{j+1}}
\]
now prove \eqref{eq:products}. The two comparisons come from
different error sums; there is no instantaneous power identity.
The full terminal angular range in Lemma~\ref{lem:graphs} gives
\eqref{eq:intervalsize}. Integrating lengths against the physical
Jacobian $s_0\,ds_0$ gives \eqref{eq:cylinderarea}, since
$\int_0^1s_0\,ds_0=1/2$. Each boundary has singleton fibre
sections and smooth endpoints for $s_0>0$, hence zero area by
Fubini. Their countable union and the vertex are null.
\end{proof}

\section{Arbitrary approximate inputs and stopped witnesses}
\label{sec:catalogue}

For any input $u$, put
$V_F(n,\delta,u)=\{x\in Q:F_u^jx\in Q^\delta,\ 0\le j\le n\}$.
The exact natural code is a proof tool, not a restriction on this set.
Define
\[
\ell_- =\min_i\frac{a_i^\beta e^{-\nu}}{a_i+\nu}>0,
\qquad C_b=\frac{2C_RC_A}{\ell_-(1-6\nu)}.
\]

\begin{lemma}[Physical first-deviation band]\label{lem:band}
Suppose an arbitrary input $u$ covers $(s_0,s_0y_0)$ through time
$n$ at tolerance $\delta$ and first differs from its natural code
at time $j\le n$. For $a=u|_{j-1}$ there is a separating point
$b_a(s_0)\in I_a(s_0)$, depending only on the common prefix and
radius, such that
\begin{equation}\label{eq:band}
|y_0-b_a(s_0)|\le C_b\delta\lambda_a^{1-\beta}/s_0.
\end{equation}
For fixed $u,j$ these states lie on one side of that boundary,
with physical area at most $C_b\delta\lambda_a^{1-\beta}$.
\end{lemma}

\begin{proof}
The common prefix is exactly feasible. On its full image graph let
$Y_*$ be the zero of $\Psi(Y)=Y-2P_0(S(Y),Y)+1$ and put
$b_a=Y_{a,s_0}^{-1}(Y_*)$. For the wrong branch $i=u_j$, direct
substitution in \eqref{eq:angularmap} gives
\begin{equation}\label{eq:wrongbranch}
|Y_j|-1=|\Psi(Y)|/P_i,\qquad
|z_j|-s_j=S(R_i/P_i)|\Psi(Y)|.
\end{equation}
At a tie both excesses are zero, so an alternate branch is included.
A necessary condition for membership in $Q^\delta$ is
$|z_j|-s_j<2\delta$, by comparison with a point of $Q$ at
Euclidean distance less than $\delta$.
Use $S\ge C_R^{-1}s_0\lambda_a^\beta$ from
\eqref{eq:products}, the global lower bound $R_i/P_i\ge\ell_-$,
and
\[
\frac{d}{dy_0}\Psi(Y_{a,s_0}(y_0))
\ge(1-6\nu)C_A^{-1}\lambda_a^{-1}.
\]
The mean value theorem proves \eqref{eq:band}. The wrong branch
fixes its side. Multiplying the width by $s_0\,ds_0$ and integrating
cancels $1/s_0$, proving the area bound without a radial cutoff.
Only the necessary condition at the first difference was used.
Later reentry cannot remove it; no cone or exact-feasibility bound
has been imposed on the subsequent approximate trajectory.
\end{proof}

Let $\mathcal L(n,t)$ stop every binary branch at its first word
with $\lambda_w\le t$, or at depth $n$ if there is no earlier
crossing. Write $N(n,t)=\#\mathcal L(n,t)$ for $0<t<1$.
This complete prefix-free counting tree is not assumed to be
an optimal catalogue. Parent/child weight addition gives
\begin{equation}\label{eq:stoppingweights}
\sum_{w\in\mathcal L(n,t)}\lambda_w=1,\qquad
\lambda_w>a_*t,\qquad N(n,t)\le(a_*t)^{-1}.
\end{equation}
Every proper prefix precedes a crossing, so the strict lower weight
bound holds also for capped leaves.

\begin{proposition}[Comparison with the actual catalogue]\label{lem:catalogue}
For fixed constants $B,C_0,C_U$, all $n\ge1$ and $0<\delta<1$,
\begin{equation}\label{eq:realcomparison}
\frac{N(n,\delta^{1/\beta})}{1+Bn}
\le r_F(n,\delta;Q)
\le N\bigl(n,(\delta/C_0)^{1/\beta}\bigr)
\le C_U\delta^{-1/\beta}.
\end{equation}
Also $r_F(n,\delta;Q)\le2^n$.
\end{proposition}

\begin{proof}
For the lower bound set $t=\delta^{1/\beta}$ and take the physical
witness $(1,y_w)$ at the midpoint of each $I_w(1)$ for the leaves
at this threshold. Their intervals partition the fibre, and
\eqref{eq:intervalsize}--\eqref{eq:stoppingweights} give successive
midpoint gaps greater than $2C_A^{-1}a_*t$.
Each midpoint is interior to every ancestor, so has natural prefix $w$.
An arbitrary input has at most one leaf as a prefix. Any other
covered witness first differs at some $j\le n$, with the one common
prefix $a=u|_{j-1}$. It is a proper prefix of the witness leaf,
so $\lambda_a>t$. Lemma~\ref{lem:band}, at radius $1$, confines
these witnesses to a one-sided band of width at most
$C_b\delta t^{1-\beta}=C_bt$. Its spacing permits at most
\[
B=\left\lceil1+\frac{C_AC_b}{2a_*}\right\rceil
\]
witnesses for one input at that time. Adding the possible leaf-prefix
witness and summing over times proves the lower bound.

For the upper bound set $t=(\delta/C_0)^{1/\beta}$ and use every
leaf as an exact prefix, followed by the same infinite tail, say
$0^\infty$. Lemma~\ref{lem:graphs} covers every radius and angle
through the leaf, including all intermediate constraints and closed
endpoints. A capped depth-$n$ leaf needs no further constrained time.
An earlier stop has radius at most $C_Rt^\beta$, hence adapted
norm at most $MC_Rt^\beta$. With
\[
C_0=2\sqrt2MC_R,\qquad C_U=C_0^{1/\beta}/a_*,
\]
the global bound \eqref{eq:globaloriginbound} keeps the entire
remaining tail within Euclidean norm $\delta/2$, strictly inside
$Q^\delta$ at every later time. The vertex is fixed.
Equation \eqref{eq:stoppingweights} gives the last inequality.
Alternatively all exactly feasible length-$n$ prefixes, extended
arbitrarily, give the $2^n$ bound.
\end{proof}

\section{Area-typical initial states and proof of the threshold}
\label{sec:threshold}

Let $c_k(x)$ denote the first $k$ letters of the natural exact code,
with the vertex assigned $0^\infty$. The continuous maps and separating
curves make every finite assigned-code set Borel measurable.

\begin{lemma}[Physical typical sets]\label{lem:typical}
For area-almost every $x\in Q$, the frequency of zero symbols
in $c_k(x)$ tends to $a_0$.
Every compact $K\subset Q$ of positive area contains a compact
positive-area set $T$ on which this convergence is uniform.
For every $0<\vartheta<h$, there is $D_\vartheta\ge1$ such that
\begin{equation}\label{eq:typicalproducts}
 D_\vartheta^{-1}e^{-(h+\vartheta)k}
 \le\lambda_{c_k(x)}
 \le D_\vartheta e^{-(h-\vartheta)k}
 \quad(x\in T,\ k\ge0).
\end{equation}
For each $\zeta\in(0,1)$ such a compact $T$ can be chosen with
$\area(T)>1-\zeta$ inside $Q$.
\end{lemma}

\begin{proof}
Normalize area on $Q$, whose area is one.
For every assigned word $w$, \eqref{eq:cylinderarea} bounds
its probability by $C_A\lambda_w$. Thus the area of states
whose length-$k$ prefix has zero frequency at distance at
least $\vartheta$ from $a_0$ is at most
\begin{equation}\label{eq:typeprob}
 C_A(k+1)e^{-2\vartheta^2 k}.
\end{equation}
Indeed, for $t=j/k$,
put $D_{\rm KL}(t\Vert a_0)=t\log(t/a_0)
+(1-t)\log((1-t)/a_1)$. Then
\[
 \binom{k}{j}a_0^j a_1^{k-j}
 \le e^{-kD_{\rm KL}(t\Vert a_0)},\qquad
 D_{\rm KL}(t\Vert a_0)\ge2(t-a_0)^2.
\]
The first inequality uses $\binom{k}{j}\le e^{kH_b(j/k)}$:
the binomial mass at $j$ with parameter $j/k$ is at most one.
The endpoint types follow by continuity. The second inequality follows from
$D_{\rm KL}''(t)=1/[t(1-t)]\ge4$, with its value and first
derivative zero at $a_0$. Summing over types gives
\eqref{eq:typeprob}. Borel--Cantelli, followed by a countable
sequence of positive $\vartheta$ tending to zero, proves almost
everywhere convergence. No independence of the state-dependent
digits has been used.

Remove the null set of finite-word boundary curves and the
vertex. Egorov's theorem and inner regularity of area give
a compact subset of any positive-area $K$ with positive area
and uniform frequency convergence. On $Q$ they give a
compact set of area greater than $1-\zeta$.
Since $-\log\lambda_{c_k(x)}/k$ is $\chi_a$ of the zero
frequency, it converges uniformly to $\chi_a(a_0)=h$ on each
chosen set. Absorb the finitely many early times into
$D_\vartheta$ to obtain \eqref{eq:typicalproducts}.
All sets are chosen before any precision exponent or sequence.
\end{proof}

\begin{lemma}[Area covered by one arbitrary input]\label{lem:typicalarea}
If $T$ satisfies \eqref{eq:typicalproducts}, then for each
$0<\vartheta<h$ there is a constant $C_\vartheta$ such that, uniformly
in $u,n,\delta$ and $0\le m\le n$,
\begin{equation}\label{eq:typicalarea}
 \area(T\cap V_F(n,\delta,u))
 \le C_\vartheta\left[e^{-(h-\vartheta)m}
       +\delta e^{(\beta-1)(h+\vartheta)m}\right].
\end{equation}
\end{lemma}

\begin{proof}
Agreement through $m$ confines the state to $\mathcal A_{u|_m}$.
If this agreement class meets $T$, its reference product is
at most $D_\vartheta e^{-(h-\vartheta)m}$, and
\eqref{eq:cylinderarea} gives the first term.

For a first disagreement at $j\le m$, a nonempty class in
$T$ has common prefix $a=u|_{j-1}=c_{j-1}(x)$ for some $x\in T$.
Equation \eqref{eq:typicalproducts} bounds its product from
below by $D_\vartheta^{-1}e^{-(h+\vartheta)(j-1)}$.
Lemma \ref{lem:band} bounds its area by
$C_bD_\vartheta^{\beta-1}\delta
 e^{(\beta-1)(h+\vartheta)(j-1)}$.
The sum over $j\le m$ is a geometric sum with ratio
greater than one, proving the second term. Empty classes
give no contribution. Later disagreement times are already
included in the agreement class through $m$.
\end{proof}

\begin{proof}[Proof of Theorem \ref{thm:main}]
The dynamical assertions follow from Lemma \ref{lem:dynamics}.
For a given positive-area compact $K$, choose $T\subset K$
as in Lemma \ref{lem:typical}. For small $\delta$ set
\[
 m=\min\left\{n,\left\lfloor
       \frac{-\log\delta}{\beta(h+\vartheta)}\right\rfloor\right\}.
\]
Then $\delta e^{\beta(h+\vartheta)m}\le1$, and
Lemma \ref{lem:typicalarea} gives
\[
 \area(T\cap V_F(n,\delta,u))\le
             2C_\vartheta e^{-(h-\vartheta)m}.
\]
Every catalogue for $K$ must cover $T$, so
\begin{equation}\label{eq:lowerK}
 r_F(n,\delta;K)\ge
       \frac{\area(T)}{2C_\vartheta}e^{(h-\vartheta)m}.
\end{equation}
For $-\log\delta_n/n\to\gamma$, passage to rates and then
$\vartheta\downarrow0$ yields the lower bound
$\min\{h,\gamma/\beta\}$. Proposition~\ref{lem:catalogue} and
$K\subset Q$ give both upper bounds in \eqref{eq:universal}.
They coincide with the lower bound when $\gamma\le\beta h$.
The argument includes $\gamma=0$ and equality at the
threshold, and needs no monotonicity of $\delta_n$.

For the high side choose once a compact $T_\zeta$ from
Lemma \ref{lem:typical} with area greater than $1-\zeta$,
and set $K_\zeta=T_\zeta\cup\{(0,0)\}$.
For every $\vartheta>0$ and all sufficiently large $n$, all
natural prefixes on $T_\zeta$ have zero frequency within
$\vartheta$ of $a_0$. The number of these prefixes is at most
\begin{equation}\label{eq:typicalcount}
 (n+1)\exp\left[n
    \max_{\substack{0\le t\le1\\|t-a_0|\le\vartheta}}H_b(t)\right].
\end{equation}
Use each as an exact length-$n$ prefix with any infinite
extension, and one extra input for the vertex. This covers
$K_\zeta$ exactly for all required times, for every
$\delta>0$. Continuity of $H_b$ gives upper rate $h$.
For every finite $\gamma>\beta h$, \eqref{eq:universal}
gives lower rate $h$ on this same set. Its rate therefore
exists and equals $h$ simultaneously for all those
exponents and all their precision sequences.

It remains to obtain a strictly larger lower rate on $Q$.
Since $a_0\ne a_1$,
\[
 \Delta=\chi_a(1/2)-h
       =(a_0-1/2)\log(a_0/a_1)>0 .
\]
For a given $\gamma>\beta h$, let
\begin{equation}\label{eq:pgamma}
 \theta_\gamma=\min\left\{\frac12,
             \frac{\gamma/\beta-h}{2\Delta}\right\},
 \qquad t_\gamma=(1-\theta_\gamma)a_0+\theta_\gamma/2 .
\end{equation}
By affinity of $\chi_a$ and strict concavity of $H_b$,
$\chi_a(t_\gamma)<\gamma/\beta$ and $H_b(t_\gamma)>h$.
Choose integers $j_n$ with $j_n/n\to t_\gamma$.
For all sufficiently large $n$, every word of length $n$
with exactly $j_n$ zeros has
\[
 -\log\lambda_w=n\chi_a(j_n/n)<-\log\delta_n/\beta.
\]
Every proper prefix has still smaller cost. All these words
are therefore terminal leaves of
$\mathcal L(n,\delta_n^{1/\beta})$.
The elementary type bound
$\binom{n}{j}\ge e^{nH_b(j/n)}/(n+1)$ and
Proposition~\ref{lem:catalogue} imply
\begin{equation}\label{eq:highlower}
 r_F(n,\delta_n;Q)\ge
   \frac{\exp[nH_b(j_n/n)]}{(n+1)(1+Bn)}.
\end{equation}
For completeness, the type bound follows since $j$ is a
mode of the binomial law with parameter $j/n$, so its
mass is at least $1/(n+1)$; endpoint types are immediate.
Taking lower rates in \eqref{eq:highlower} yields
$H_b(t_\gamma)>h$, proving \eqref{eq:high}.
The strict gap may tend to zero as $\gamma\downarrow\beta h$.
\end{proof}

\begin{corollary}[Matched linear systems]\label{cor:linear}
For
\[
 A_0(s,z)=(a_0^\beta s,a_0^{\beta-1}(z+a_1s)),\qquad
 A_1(s,z)=(a_1^\beta s,a_1^{\beta-1}(z-a_0s)),
\]
the conclusions of Theorem \ref{thm:main} hold with the
same threshold $\beta h$.
\end{corollary}
\begin{proof}
Set $\eta=\xi_0=\xi_1=0$ in Theorem~\ref{thm:main}.
The independent radial and angular constants are then $C_R=C_A=1$.
\end{proof}


\section{Radius-autonomous matching and mismatch}\label{sec:radius}

We now use \eqref{eq:G}, with parameters \eqref{eq:radialepsilon}.
The quantities in this section are independent of the feedback
constants. For words let
$\lambda_w=\prod_j a_{w_j}$ and $R_w=\prod_j\rho_{w_j}$.
Write $\beta_{\min}=\min_i\beta_i$ and put
\begin{gather*}
q_-=a_*-\ee/2,\qquad q_+=1-q_-,\qquad r_R=q_+^{\beta_{\min}}<1,\\
C_R^P=\exp\!\left(\frac{\ee}{q_-(1-r_R)}\right),\qquad
C_R^s=(C_R^P)^{\beta_{\max}}.
\end{gather*}

\begin{lemma}[Actual strips and radial products]\label{lem:radialstrips}
The maps \eqref{eq:G} are global smooth diffeomorphisms and admit
a common origin-contracting norm on the whole plane. At a fixed
initial radius, each complete feasible word has a closed angular
interval $I_w(s)$, and every complete prefix-free tree partitions
that fibre with endpoints included. Along a word set
\[
\mathcal Q_w(s)=\prod_jq_{w_{j+1}}(s_j),\qquad
\mathcal R_w(s)=\prod_jq_{w_{j+1}}(s_j)^{\beta_{w_{j+1}}}.
\]
Then, uniformly in word and radius $0\le s\le1$,
\begin{equation}\label{eq:radialproducts}
(C_R^P)^{-1}\lambda_w\le\mathcal Q_w\le C_R^P\lambda_w,
\qquad (C_R^s)^{-1}R_w\le\mathcal R_w\le C_R^sR_w,
\end{equation}
$s_{|w|}=s\mathcal R_w(s)$ and $|I_w(s)|=2\mathcal Q_w(s)$.
Physical strip area is between $(C_R^P)^{-1}\lambda_w$ and
$C_R^P\lambda_w$.
\end{lemma}

\begin{proof}
For $f(s)=s/(1+s^2)$, $|f|\le1/2$ and $|sf'|\le1/2$.
Thus $g_i(s)=sq_i(s)^{\beta_i}$ has derivative
\[
g_i'(s)=q_i^{\beta_i-1}(q_i+\beta_i s q_i')>0,
\]
with the bracket at least $a_*-(\beta_{\max}+1)\ee/2>0$.
Since $q_i(s)\to a_i$ at infinity, $g_i$ is onto $\R$.
Its smooth inverse recovers $s$, then the positive coefficient
$q_i^{\beta_i-1}$ recovers $z$. This proves global invertibility.
With $M_R\ge1$, $M_R>r_R/(1-q_+^{\beta_{\min}-1})$,
the norm $\max(M_R|s|,|z|)$ contracts relative to the origin
by at most
$\max\{r_R,q_+^{\beta_{\min}-1}+r_R/M_R\}<1$.

At radius $s$, the angular maps are
$T_{0,s}(y)=(y+q_1)/q_0$ and $T_{1,s}(y)=(y-q_0)/q_1$.
The full domains are $[-1,\xi(s)]$, $[\xi(s),1]$, where
$\xi=q_0-q_1$, with inverse branches
$B_{0,s}(y)=q_0(s)y-q_1(s)$ and
$B_{1,s}(y)=q_1(s)y+q_0(s)$.
The radius after the first letter does not depend on the angle, so
\begin{equation}\label{eq:radialintervals}
I_{iv}(s)=B_{i,s}(I_v(g_i(s))).
\end{equation}
This induction retains all intermediate constraints and gives the
partition and length assertions. Controlled feasibility follows.

The radial path obeys $s_j\le r_R^js$. Since
$|q_i(s_j)-a_i|\le\ee s_j$, summing logarithmic errors proves
\eqref{eq:radialproducts}; the radial error is at most
$\beta_{\max}$ times the angular one. Integration against $s\,ds$
gives the area bounds. Finite-word boundary curves have zero area,
and the vertex is fixed.
\end{proof}

\begin{lemma}[Relative radial-stop witness comparison]\label{lem:radialcomparison}
Let $\mathcal L_R(n,t)$ stop words at first $R_w\le t$, or at
depth $n$, and let $N_R(n,t)$ be its size. Fixed constants
$B_R,C_{0,R}$ satisfy
\begin{equation}\label{eq:radialcomparison}
\frac{N_R(n,\delta)}{1+B_Rn}
\le r_G(n,\delta;Q)
\le N_R(n,\delta/C_{0,R})
\end{equation}
for all horizons and sufficiently small positive tolerances.
For a first difference at time $j$, common prefix $a$, an arbitrary
covered physical state lies in a one-sided band of width
\begin{equation}\label{eq:radialband}
|y-b_a(s)|\le C_{b,R}\delta\lambda_a/(sR_a),\qquad
C_{b,R}=2C_R^P C_R^s\max_i(a_i/\rho_i).
\end{equation}
Its physical area is at most $C_{b,R}\delta\lambda_a/R_a$.
\end{lemma}

\begin{proof}
The common word is exact. Its angle map is affine with slope
$\mathcal Q_a(s)^{-1}$. For the differing branch $i$, the excess
at the first difference is
\[
|y_j|-1=|y-b_a(s)|/\mathcal Q_{ai}(s),\qquad
s_j=s\mathcal R_{ai}(s).
\]
The necessary condition $|z_j|-s_j<2\delta$ and
\eqref{eq:radialproducts} give \eqref{eq:radialband}.
The side is fixed by the differing input; integration cancels $1/s$.
This includes ties and subsequent reentry.

For the lower bound use midpoint witnesses at radius $1$ for
leaves stopped at $t=\delta$. Write $\rho_* =\min_i\rho_i$.
Every stopped or capped leaf satisfies $R_w>\rho_*t$.
For any descendant $w=av$ of a fixed proper prefix $a$,
$\lambda_v\ge R_v=R_w/R_a>\rho_*t/R_a$, because $\beta_i>1$.
By Lemma~\ref{lem:radialstrips}, all leaf intervals below $a$
have length greater than
$2(C_R^P)^{-1}\lambda_a\rho_*t/R_a$.
Their ordered midpoints have at least that spacing.
For one input and one first-difference time, \eqref{eq:radialband}
at $s=1$ confines these descendant witnesses to a band of width
$C_{b,R}\delta\lambda_a/R_a$. Hence at most
\[
B_R=\left\lceil1+\frac{C_{b,R}C_R^P}{2\rho_*}\right\rceil
\]
can be covered at that time. Add the one possible leaf prefix
of the input and sum over times to get $1+B_Rn$.
This is a conditional spacing comparison within the shared
prefix, rather than a global minimum-length assertion.

For the upper bound, the full leaf intervals of
\eqref{eq:radialintervals} cover each physical fibre exactly
through the leaf. Append the same infinite tail to each leaf
at $t=\delta/C_{0,R}$, with
$C_{0,R}=2\sqrt2M_R C_R^s$. An early terminal radius is
at most $C_R^st$, and the global origin-contraction norm of
Lemma~\ref{lem:radialstrips} keeps the entire tail within
Euclidean distance $\delta/2$ of the origin. Capped leaves
are exactly feasible through $n$. All endpoints and the vertex
are covered, proving \eqref{eq:radialcomparison}.
\end{proof}

\begin{proof}[Proof of Theorem~\ref{thm:mismatch}]
Lemma~\ref{lem:radialstrips} proves the dynamical assertions.
Write $V_G(n,\delta,u)=\{x\in Q:G_u^jx\in Q^\delta,
0\le j\le n\}$ for the actual all-time initial set.
For $0<\gamma<\gamma_0$, the bound $R_w\le e^{-\gamma_0n}$
at depth $n$ implies that every branch crosses either of the
thresholds in \eqref{eq:radialcomparison} by time $n$, eventually
for any precision sequence of exponent $\gamma$.
For an uncapped radial stopping tree, replacement of a parent
weight $R_w^D$ by its children preserves total weight, since
$\rho_0^D+\rho_1^D=1$. All leaves have
$\rho_*t<R_w\le t$, and thus
\begin{equation}\label{eq:radialcount}
t^{-D}\le N_R(t)\le(\rho_*t)^{-D}.
\end{equation}
The comparison proves the full-set limit $\gamma D$.

By the area bound of Lemma~\ref{lem:radialstrips}, the type
probabilities of actual assigned codes are bounded by $C_R^P$
times their Bernoulli weights. The proof of Lemma~\ref{lem:typical}
therefore gives one uniformly typical compact $T_\zeta$ of area
greater than $1-\zeta$. Set $K_\zeta=T_\zeta\cup\{0\}$ before
choosing any exponent, sequence, or horizon. On this same set,
for each $\vartheta>0$ a constant $A_\vartheta\ge1$ satisfies
\begin{align}
A_\vartheta^{-1}e^{-(h+\vartheta)k}&\le\lambda_{c_k(x)}
\le A_\vartheta e^{-(h-\vartheta)k},\nonumber\\
A_\vartheta^{-1}e^{-(\chi+\vartheta)k}&\le R_{c_k(x)}
\le A_\vartheta e^{-(\chi-\vartheta)k}.\label{eq:twotypical}
\end{align}
Both follow from the same uniform frequency convergence.
For any positive-area compact $K$, the same construction supplies
a positive-area compact $T\subset K$ with these bounds.

Agreement with an arbitrary input through $m$ has area at most
$C_R^P A_\vartheta e^{-(h-\vartheta)m}$. A nonempty first-difference
class at time $j\le m$ has a common typical prefix. Equations
\eqref{eq:radialband} and \eqref{eq:twotypical} bound its
area by a constant times
$\delta e^{(\chi-h+2\vartheta)(j-1)}$.
Since $\chi>h$, summation gives, uniformly in the input,
\begin{equation}\label{eq:radialtypicalarea}
\area(T\cap V_G(n,\delta,u))\le
A'_\vartheta[e^{-(h-\vartheta)m}+\delta e^{(\chi-h+2\vartheta)m}].
\end{equation}
For $m=\min\{n,\lfloor(-\log\delta)/(\chi+\vartheta)\rfloor\}$,
the second term is at most $e^{-(h-\vartheta)m}$.
Since $\gamma<\gamma_0\le\chi$, covering the positive area
of $T$ yields lower rate $\gamma h/\chi$, after $\vartheta\downarrow0$.

For the matching upper rate on $K_\zeta$, use all radial
stopping prefixes realized on $T_\zeta$, with the full tails
from Lemma~\ref{lem:radialcomparison}. Put $T=-\log\delta$
and $c_{\max}=-\log\rho_*$. Each such stop has length at most
\[
\ell\le\frac{T+\log C_{0,R}+c_{\max}+\log A_\vartheta}{\chi-\vartheta}.
\]
At a fixed length $\ell$, distinct realized words have
$\lambda_w\ge A_\vartheta^{-1}e^{-(h+\vartheta)\ell}$, whereas
all length-$\ell$ word weights sum to one. Their number is
at most $A_\vartheta e^{(h+\vartheta)\ell}$. Summing up to the displayed
length, and adding one input for the apex, gives upper rate
$\gamma(h+\vartheta)/(\chi-\vartheta)$. Taking $\vartheta\downarrow0$
proves the second limit in \eqref{eq:mismatchrates}, on one
fixed $K_\zeta$ for all the stated sequences.

Finally let $b_i=\rho_i^D$, so $b_0+b_1=1$. Then
\begin{equation}\label{eq:KLmismatch}
D\chi-h=\sum_i a_i\log(a_i/b_i)\ge0.
\end{equation}
The inequality follows from $-\log t\ge1-t$; equality holds
only when $a_i=b_i$ for both digits. This would imply
$\beta_0D=\beta_1D=1$, which is impossible under mismatch.
Hence $D>h/\chi$. The root and the strict entropy comparison
are classical counting, used only after the physical catalogue
comparison has been established.
\end{proof}

\begin{proposition}[Matched radius-autonomous subfamily]\label{prop:radialmatched}
If $\beta_0=\beta_1=\beta$, the system \eqref{eq:G}, on its
full range \eqref{eq:radialepsilon}, has all the rate conclusions
of Theorem~\ref{thm:main}, including the single compact set for
every finite high-side exponent.
\end{proposition}

\begin{proof}
Here $R_w=\lambda_w^\beta$ and $\mathcal R_w=\mathcal Q_w^\beta$.
Lemma~\ref{lem:radialstrips} gives the full intervals and area
bounds required in Lemma~\ref{lem:typical}.
Equation \eqref{eq:radialband} becomes a one-sided physical
band of area at most $C_{b,R}\delta\lambda_a^{1-\beta}$.
Equation \eqref{eq:radialcomparison} becomes
\eqref{eq:realcomparison} with fixed constants and angular
stopping thresholds $\delta^{1/\beta}$ and
$(\delta/C_{0,R})^{1/\beta}$. Its upper count follows from
\eqref{eq:stoppingweights}; the $2^n$ exact bound also holds.
Thus the proofs of Lemma~\ref{lem:typicalarea} and
Theorem~\ref{thm:main} use exactly these already derived
estimates: choose the same $m$ for the lower bound, one
uniformly typical compact before all exponents for the high-side
upper bound, and the type $t_\gamma$ in \eqref{eq:pgamma}
for the full-set strict lower bound. Every inequality and tail
needed for those deductions has been supplied in this section.
\end{proof}

\begin{corollary}[Matching criterion within the radius-autonomous class]\label{cor:criterion}
For a fixed system \eqref{eq:G}, there is a nonempty interval
of exponents adjacent to zero on which all fixed positive-area
compact initial sets have the same rate for every corresponding
precision sequence if and only if $\beta_0=\beta_1$.
\end{corollary}
\begin{proof}
Under equality, Proposition~\ref{prop:radialmatched} gives
the interval $0<\gamma\le\beta h$. Under inequality,
Theorem~\ref{thm:mismatch} separates $Q$ and one fixed
nearly full-area compact set for every $0<\gamma<\gamma_0$.
\end{proof}

This criterion concerns only \eqref{eq:G}. It does not give
a necessity statement for \eqref{eq:F}, overlapping feasible
branches, or general smooth perturbations.


\section{Independent factors beyond coordinate reduction}
\label{sec:conjugacy}

Some perturbations, including the zero triple, reduce to old forms.
The next construction excludes this explanation for genuine instances
of \eqref{eq:F}. For comparison define the exactly matched rational
feedback family
\begin{gather}
0\le\ee<\frac{a_*}{16(\beta+1)},\qquad
\widehat p_0(s,z)=a_0+\frac{\ee z}{1+s^2+z^2},\quad
\widehat p_1(s,z)=a_1-\frac{\ee z}{1+s^2+z^2},\label{eq:feedbackepsilon}\\
H_0=(s\widehat p_0^\beta,\widehat p_0^{\beta-1}(z+\widehat p_1s)),
\quad
H_1=(s\widehat p_1^\beta,\widehat p_1^{\beta-1}(z-\widehat p_0s)).
\label{eq:H}
\end{gather}
This family is not asserted to satisfy the compact-support
hypothesis of Theorem~\ref{thm:main}. Nor is every system in
\eqref{eq:F} asserted to reduce to \eqref{eq:H}.

\begin{proposition}[A feasible-image obstruction]\label{prop:conjugacy}
There are systems in the full class of Theorem~\ref{thm:main}
with $R_i\ne p_i^\beta$ which admit no simultaneous ambient
homeomorphism preserving $Q$ and conjugating both controls to any
system \eqref{eq:H} or \eqref{eq:G}, including the linear endpoints
and allowing label exchange. The identities are required on $Q$,
on a domain containing $Q$ and its needed one-step images.
In particular no $C^1$ or bi-Lipschitz coordinate change of this
kind works for these instances.
\end{proposition}

\begin{proof}
Let $E(v)=e^{-1/v}$ for $v>0$ and $E(v)=0$ otherwise, and put
\[
\kappa(x)=\frac{E(49/16-|x|^2)}
 {E(49/16-|x|^2)+E(|x|^2-9/4)}.
\]
Its denominator is positive. The function is one on the closed
ball of radius $3/2$ and zero outside radius $7/4$, hence is
smooth and compactly supported in $B(0,2)$.
Define
\begin{gather*}
f(s,z)=\kappa(s,z)\sin(4\pi z),\quad
C_f=\max\{1,\max_{|\alpha|\le2}\|\partial^\alpha f\|_\infty\},\\
\bar\nu=1/4000,\quad u=\bar\nu/(2C_f),\quad d_1=u/64,\quad
\beta=2,\quad a_0=(1+d_1)/2,\quad a_1=(1-d_1)/2.
\end{gather*}
Take $\eta=0$, $\xi_0=uf$, and $\xi_1=-uf$.
They vanish at zero and satisfy \eqref{eq:functionbounds} with
bound at most $\bar\nu/2$. Here $a_*>0.49$ and $a^+<0.51$.
The first term in \eqref{eq:smallthreshold} exceeds $\bar\nu$;
so does the second, since $-\log0.51>2/5$ and
$4(1+2/0.49)<24$. Thus this construction belongs to the theorem
with $\nu=\bar\nu$.

Let $D_i=Q\cap F_i^{-1}Q$ and $J_i=F_i(D_i)$.
The factors $p_i=a_i$ are constant, so the source angles at a
terminal angle $Y$ are
$y_0(Y)=a_0Y-a_1$ and $y_1(Y)=a_1Y+a_0$.
Every source in $Q$ lies where $\kappa=1$.
For fixed $Y$, its output radius is
\[
a_i^2s\exp[(-1)^iu\sin(4\pi sy_i(Y))].
\]
It increases strictly with $s\in[0,1]$, because the derivative
bracket is at least $1-4\pi u>0$.
Thus $J_i$ is exactly the radial region $0\le S\le U_i(Y)$,
$-1\le Y\le1$, where
\[
U_0(Y)=a_0^2e^{u\sin(4\pi y_0(Y))},\qquad
U_1(Y)=a_1^2e^{-u\sin(4\pi y_1(Y))}.
\]
The parts of their boundaries inside $Q$ are precisely the upper
radial arcs, for $-1<Y<1$; both lie below the cap $s=1$.
Their log-radius difference is
\begin{equation}\label{eq:intersectiondifference}
\mathcal D(Y)=2\log\frac{1+d_1}{1-d_1}
 +2u\sin(2\pi(Y+d_1))\cos(2\pi d_1Y).
\end{equation}
Its constant term is at most $8d_1=u/8$.
At $Y=-3/4,-1/4,1/4,3/4$ the values of $\sin(2\pi Y)$ have
alternating signs $+,-,+,-$ and magnitude one.
The error in replacing the sine/cosine factor in
\eqref{eq:intersectiondifference} by this sine is at most
$2\pi d_1+2\pi^2d_1^2<1/4$.
Therefore $\mathcal D$ has the same alternating signs and at
least three distinct interior zeros. It is analytic on a
neighbourhood of $[-1,1]$ and is not identically zero, so it has
only finitely many zeros there. Consequently
\begin{equation}\label{eq:boundaryinvariant}
\partial J_0\cap\partial J_1\cap\operatorname{int}Q
\end{equation}
is a finite set with at least three points.

For a system \eqref{eq:G}, the radial maps are increasing and
independent of the angle, by Lemma~\ref{lem:radialstrips}.
Its feasible images are triangles with cap radii $g_i(1)$.
Their interior boundary arcs are either disjoint or coincide
in an entire arc. Neither gives the intersection in
\eqref{eq:boundaryinvariant}.

For \eqref{eq:H}, put $\widehat P_i(s,y)=\widehat p_i(s,sy)$ and
$\widehat p_-=a_*-\ee/2$. The one-step angle maps are increasing
on each feasible domain. Indeed
$|\partial_y\widehat P_i|\le\ee s$ gives
$\partial_yY\ge(1-2\ee/\widehat p_-)/\widehat P_i>0$.
The determinant in physical coordinates is
\[
\widehat p_i^{2\beta-2}
 [\widehat p_i+\beta s(\widehat p_i)_s+
 ((\beta-1)z+\sigma_i s)(\widehat p_i)_z]>0.
\]
To check positivity, for $f=z/(1+s^2+z^2)$ use
$|sf_s|\le1$, $|zf_z|\le1/2$, and $|sf_z|\le1/2$;
the derivative contribution is at most $3\beta\ee/2<\widehat p_-$.
At fixed output angle, the output radius therefore increases
with source radius: in $(s,y)$ coordinates its derivative is
the positive Jacobian divided by $\partial_yY$.
For every $s$ the feasible angular maps have full range $[-1,1]$.
The interior upper boundaries of the two images consequently
come from the source cap $s=1$.

On that cap $\widehat P_0=a_0+\ee y/(2+y^2)$ is increasing in
$y$, while $\widehat P_1=a_1-\ee y/(2+y^2)$ is decreasing,
since $(y/(2+y^2))'=(2-y^2)/(2+y^2)^2>0$.
The cap angle maps are increasing, and the cap radii are the
same positive powers of these factors. For $\ee>0$, the
two upper arcs are therefore respectively increasing and
decreasing as functions of $Y$, and meet at most once.
For $\ee=0$ their constant radii are distinct, since
$a_0\ne a_1$.

A simultaneous ambient homeomorphism taking $Q$ onto $Q$ maps
each feasible domain and image to its target counterpart.
It maps boundaries and the interior of $Q$ bijectively, preserving
the cardinality in \eqref{eq:boundaryinvariant}; label exchange
does not change it. The finite intersection with at least three
points contradicts both target possibilities.
Finally $p_i^\beta=a_i^2$, whereas $R_i=a_i^2e^{\pm uf}$ is
nonconstant on $Q$, verifying the failure of instantaneous matching.
\end{proof}

The proposition is an existence statement for the displayed instances,
not a nonconjugacy assertion for every triple in the main theorem.
It excludes only simultaneous changes preserving $Q$ and the actual
control identities. It makes no claim about changes that alter the
constraint, conjugate only one control, or target an unspecified
larger model class.

\section{Local matching does not suffice}
\label{sec:local}

The following example keeps the common attracting origin
and even agrees with the linear maps on a neighbourhood
of it. It violates global injectivity by collapsing
positive-area regions onto an invariant ray.

\begin{proposition}\label{prop:collapse}
For every choice of $a_0,a_1,\beta$ in
\eqref{eq:parameters} and every $0<\sigma<1/2$, there are
$C^\infty$ maps $S_0,S_1:\R^2\to\R^2$ with these properties:
\begin{enumerate}
 \item They equal $A_0,A_1$ on $\{s\le\sigma\}$, so
       all derivatives at the origin agree.
 \item All inputs drive bounded sets uniformly
       exponentially to the origin, and $Q$ is
       controlled invariant.
 \item The compact set
       $K_\sigma=Q\cap\{s\ge2\sigma\}$ has area
       $1-4\sigma^2$ and $r_S(n,\delta;K_\sigma)=1$
       for every $n$ and every $\delta>0$.
 \item For every finite $\gamma\ge0$ and every sequence
       $\delta_n\to0$ with $-\log\delta_n/n\to\gamma$,
       the full-set cost satisfies
       \[
       \min\{h,\gamma/\beta\}
       \le\liminf\frac{\log r_S(n,\delta_n;Q)}n
       \le\limsup\frac{\log r_S(n,\delta_n;Q)}n
       \le\min\{\log2,\gamma/\beta\}.
       \]
\end{enumerate}
In particular, for $0<\gamma\le\beta h$ the full-set rate
is $\gamma/\beta$, but the positive-area set $K_\sigma$
has rate zero.
\end{proposition}

\begin{proof}
Let $e(t)=0$ for $t\le0$ and $e(t)=\exp(-1/t)$ for $t>0$,
and define
\[
 \psi(s)=\frac{e(s-\sigma)}
               {e(s-\sigma)+e(2\sigma-s)}.
\]
The denominator is positive everywhere. Thus $\psi$ is
smooth, zero for $s\le\sigma$, and one for $s\ge2\sigma$.
Put $\rho_i=a_i^\beta$ and define
\begin{align}
 S_0(s,z)&=\bigl(\rho_0s,\
 (1-\psi(s))a_0^{\beta-1}(z+a_1s)-\psi(s)\rho_0s\bigr),
       \label{eq:S0}\\
 S_1(s,z)&=\bigl(\rho_1s,\
 (1-\psi(s))a_1^{\beta-1}(z-a_0s)+\psi(s)\rho_1s\bigr).
       \label{eq:S1}
\end{align}
They have the claimed local identity.
Each $S_i(s,z)$ is a convex combination of $A_i(s,z)$
and $R_i(s,z)$, where
$R_0=(\rho_0s,-\rho_0s)$ and
$R_1=(\rho_1s,\rho_1s)$.
All four linear maps contract toward zero in one adapted
norm $\max\{M|s|,|z|\}$ for sufficiently large $M$:
the second-coordinate factor for $A_i$ is at most
$a_i^{\beta-1}+a_i^{\beta-1}a_{1-i}/M<1$, and
$R_i$ has factor at most $\rho_i$.
Convexity of the norm gives the same contraction bound
for $S_i$ on all of $\R^2$.

For controlled feasibility let $t=\psi(s)$, $y=z/s$,
and $\xi=a_0-a_1$. If $t<1$ the full feasible intervals
are
\begin{align*}
 J_0(s)&=\left[-1,\min\left\{1,
                         \xi+\frac{2a_0t}{1-t}\right\}\right],\\
 J_1(s)&=\left[\max\left\{-1,
                         \xi-\frac{2a_1t}{1-t}\right\},1\right].
\end{align*}
They follow by imposing $-1\le z'/s'\le1$ in
\eqref{eq:S0}--\eqref{eq:S1}, and their union is
$[-1,1]$. For $t=1$ both branches are feasible on the
whole fibre. The radius always decreases, so $Q$ is
controlled invariant.

On $K_\sigma$, one application of $S_0$ sends every
state to the lower ray. At every radius
$S_0(s,-s)=(\rho_0s,-\rho_0s)$, so input $0^\infty$
keeps the whole set exactly in $Q$. Its area is
$\int_{2\sigma}^1 2s\,ds=1-4\sigma^2$.

To obtain the full-set lower bound, use
$Q_c=Q\cap\{s\le c\}$ with $c=\sigma/2$.
For every input, the radius of a state in $Q_c$ stays
at most $c$. Thus $\psi=0$ forever along that trajectory,
even when its angle leaves $[-1,1]$, and its entire
trajectory is exactly that of $A_i$.
Consequently
$r_S(n,\delta;Q)\ge r_A(n,\delta;Q_c)$.
Corollary \ref{cor:linear} and \eqref{eq:universal}
give the lower bound.

For the upper bound choose an integer $t_0$ such that
$(\max_i\rho_i)^{t_0}\le\sigma$.
Every initial state admits an exact feasible prefix of
length $t_0$, ending in the region $s\le\sigma$.
All subsequent trajectories there coincide with the
linear system, for all inputs, since the radius continues
to decrease independently of the angle.
Concatenate all $2^{t_0}$ possible prefixes with a
catalogue for the linear full set over the remaining
horizon. For $n\ge t_0$ this gives
\[
 r_S(n,\delta;Q)\le
       2^{t_0}r_A(n-t_0,\delta;Q)
       \le 2^{t_0}C_U\delta^{-1/\beta}.
\]
The $2^n$ exact-prefix bound also holds.
These yield the upper rates for every stated precision
sequence. Finally
$\det DS_i=(1-\psi(s))a_i^{2\beta-1}$ vanishes on the
saturated region, identifying the noninjective mechanism.
\end{proof}

For fixed $\sigma$ this proposition gives one large
zero-cost initial set in one fixed system. Making its
area arbitrarily close to one changes $\sigma$ and the
system. This differs from Theorem \ref{thm:main}, where
all $K_\zeta$ are chosen inside a single fixed system.
The collapse example is a counterexample to sufficiency
of local matching and uniform convergence; it is not a
counterexample to the reversible theorem.



\section{Scope and consequences}

The common-cost boundary concerns coverage of actual initial states.
Below it, a positive-area typical portion gives the lower bound for
every compact initial set, including sets with empty interior.
Above it, rare terminal types force additional controls for $Q$,
while one uniformly typical compact set excludes those deviations.
The choice of that compact set precedes every precision exponent.

The main theorem removes the instantaneous relation between branch
lengths and radial factors. Origin contraction makes the two error
sums finite, but their comparison to the real catalogue still requires
the complete feasible graphs and the physical first-deviation band.
The $1+Bn$ witness bound is what prevents arbitrary approximate
controls, including those that reenter $Q$, from saving an
exponential factor over the stopping count.

The mismatch theorem gives a different comparison. Its two radial
first derivatives have unequal powers of the limiting branch
proportions. Relative-prefix spacing connects the classical root
$D$ to the actual full-set rate, and uniformly typical physical
sets attain the smaller rate. Necessity of equal powers is proved
only in that radius-autonomous class. Neither theorem provides a
general necessary and sufficient condition for feedback systems.

The perturbation class retains fixed compact support, small
derivative bounds, matched first jets, and complementary feasible
branches without interior overlap. These are global conditions:
local matching and convergence alone fail by Section~\ref{sec:local}.
Neither theorem gives a general classification or the full-set
high-precision spectrum and its limit.

\subsection*{Generative AI disclosure}
Generative AI was used extensively in developing mathematical
arguments, checking proof chains and related literature, drafting
and restructuring the English text, and producing the LaTeX
source. These activities do not constitute verification by a
human author. Author identity, affiliations, authorship approval
and the scope of human mathematical review remain to be supplied
and confirmed before any submission.

\begin{thebibliography}{99}
\raggedright

\bibitem{CZ2024}
Z. Chen and X. Zhong,
Invariance complexity and equi-invariability for control systems,
\emph{J. Math. Anal. Appl.}, \textbf{539} (2024), 128533.
\doi{10.1016/j.jmaa.2024.128533}.

\bibitem{CHZ2026}
H. Chen, Y. Huang and X. Zhong,
Variational principles of invariance entropy dimension,
\emph{J. Differential Equations}, \textbf{453}, Part 1 (2026), 113819.
\doi{10.1016/j.jde.2025.113819}.

\bibitem{Colonius2012}
F. Colonius,
Minimal bit rates and entropy for exponential stabilization,
\emph{SIAM J. Control Optim.}, \textbf{50} (2012), 2988--3010.
\doi{10.1137/110829271}.

\bibitem{CH2021}
F. Colonius and B. Hamzi,
Entropy for practical stabilization,
\emph{SIAM J. Control Optim.}, \textbf{59} (2021), 2195--2222.
\doi{10.1137/20M1367775}.

\bibitem{Csiszar1998}
I. Csisz\'ar,
The method of types,
\emph{IEEE Trans. Inform. Theory}, \textbf{44} (1998), 2505--2523.
\doi{10.1109/18.720546}.

\bibitem{Hutchinson1981}
J. E. Hutchinson,
Fractals and self-similarity,
\emph{Indiana Univ. Math. J.}, \textbf{30} (1981), 713--747.
\doi{10.1512/iumj.1981.30.30055}.


\bibitem{KD2018}
C. Kawan and A. Da Silva,
Invariance entropy for a class of partially hyperbolic sets,
\emph{Math. Control Signals Syst.}, \textbf{30} (2018), Article 18.
\doi{10.1007/s00498-018-0224-2}.

\bibitem{Kawan2013}
C. Kawan,
\emph{Invariance Entropy for Deterministic Control Systems:
An Introduction}, Lecture Notes in Mathematics, 2089,
Springer, 2013.
\doi{10.1007/978-3-319-01288-9}.

\bibitem{NWH2022}
X. Nie, T. Wang and Y. Huang,
Measure-theoretic invariance entropy and variational principles
for control systems,
\emph{J. Differential Equations}, \textbf{321} (2022), 318--348.
\doi{10.1016/j.jde.2022.03.013}.

\bibitem{WHS2019}
T. Wang, Y. Huang and H.-W. Sun,
Measure-theoretic invariance entropy for control systems,
\emph{SIAM J. Control Optim.}, \textbf{57} (2019), 310--333.
\doi{10.1137/18M1197862}.

\bibitem{YCYZ2025}
R. Yang, E. Chen, J. Yang and X. Zhou,
Bowen's equations for invariance pressure of control systems,
\emph{SIAM J. Control Optim.}, \textbf{63} (2025), 1104--1128.
\doi{10.1137/23M1607684}.

\end{thebibliography}
\end{document}
